\documentclass[10pt,a4paper]{amsart}
\usepackage[margin=19mm]{geometry}
\usepackage{amsmath,amssymb,mathtools}
\usepackage[scr=rsfs,cal=euler]{mathalfa}
\usepackage{xfrac}
\usepackage[unicode,hidelinks,bookmarks=false]{hyperref}
\newcommand{\ie}{i.e.\ }
\newcommand{\cf}{cf.\ }
\newcommand{\resp}{respectively\ }
\newcommand{\Subsection}{\subsection}
\providecommand{\nobreakdash}{\nobreak\hbox{-}}
\setlength{\emergencystretch}{2em}
\multlinegap0pt
\usepackage{bm}
\usepackage{amssymb}
\usepackage{mathtools}
\usepackage{xcolor}
\newtheorem{prop}{Prop.}
\title{Framing structural identifiability in terms of parameter symmetries}
\author{Johannes G Borgqvist, Alexander P Browning, Fredrik Ohlsson, Ruth E Baker}
\date{Source: arXiv:2410.03757v3 (CC BY 4.0)}
\begin{document}
\maketitle

\noindent\emph{Source and licence.} Framing structural identifiability in terms of parameter symmetries. Version arXiv:2410.03757v3 (CC BY 4.0), distributed by arXiv under the Creative Commons Attribution 4.0 International licence, http://creativecommons.org/licenses/by/4.0/, as recorded in the arXiv licence metadata for this exact version and shown on the abstract page https://arxiv.org/abs/2410.03757v3. Full licence notice: \texttt{licenses/arxiv-2410.03757-CC-BY-4.0.txt}.\par\smallskip

\noindent\textit{Editorial scope.} Counted unit: the article's prop prop:inv (source0.tex lines 373--376). The article's complete original proof of it is delivered with it (source0.tex lines 377--391, one \texttt{proof} environment of 15 lines). No mathematics is written, repaired or re-derived: the statement and the proof are the article's own lines, in the article's own notation, and no retained line is substituted or re-flowed.  No further source-local context is needed: every cross-reference of every retained line resolves inside the two delivered spans.  Nothing was omitted from any retained span, so the package carries no omission record.  No retained line contains a citation, so the wrapper carries no bibliography and no bibliography entry.  No retained line contains a figure environment, an image-inclusion command or drawing code, so no image is delivered and the package declares no asset dependency.  Editorial formatting only, in addition: an AMS \texttt{amsart} preamble that restates the article's own declarations and macros used by the retained lines, namely the environments \texttt{prop} on separate counters, together with the standard packages those lines need.  The retained environments print in this document as Proposition 1 in order of appearance, whereas the article prints the same items with the numbering of its own full document.  The complete native source of the article as distributed in the arXiv e-print for this exact version is bundled unchanged as \texttt{sources/166316/source0.tex}.
\begin{prop}[Functions of differential invariants are themselves differential invariants]
Let $X_{\boldsymbol{\theta}}$ generate a parameter symmetry with $p$ parameters and denote its functionally independent parameter invariants by $I_{1},\ldots,I_{\tilde{p}}$ where the number of invariants is an integer $\tilde{p}\in\{2,\ldots,p\}$. Then any differentiable function $F\left(I_{1},\ldots,I_{\tilde{p}}\right)$ of these invariants is itself a differential invariant.
\label{prop:inv}
\end{prop}

\begin{proof}
By definition, we have that $X_{\boldsymbol{\theta}}(I_{j})=0$ $\forall j\in\left\{1,\ldots,\tilde{p}\right\}$ and we need to show that $X_{\boldsymbol{\theta}}\left(F\left(I_{1},\ldots,I_{\tilde{p}}\right)\right)=0$. By the chain rule, it follows that
\begin{equation}
\frac{\partial F}{\partial\theta_{\ell}}=\frac{\partial}{\partial\theta_{\ell}}\left(F\left(I_{1},\ldots,I_{\tilde{p}}\right)\right)=\sum_{j=1}^{\tilde{p}}\frac{\partial I_{j}}{\partial\theta_{\ell}}\frac{\partial F}{\partial I_{j}},\quad\ell\in\left\{1,\ldots,p\right\}.
\end{equation}
Thus, we have
\begin{align}
X_{\boldsymbol{\theta}}\left(F\left(I_{1},\ldots,I_{\tilde{p}}\right)\right)&=\sum_{\ell=1}^{p}\chi_{\ell}\frac{\partial F}{\partial\theta_{\ell}}\nonumber\\
&=\sum_{\ell=1}^{p}\chi_{\ell}\sum_{j=1}^{\tilde{p}}\frac{\partial I_{j}}{\partial\theta_{\ell}}\frac{\partial F}{\partial I_{j}}\nonumber\\
&=\sum_{\ell=1}^{p}\sum_{j=1}^{\tilde{p}}\chi_{\ell}\frac{\partial I_{j}}{\partial\theta_{\ell}}\frac{\partial F}{\partial I_{j}}\nonumber\\
    &=\sum_{j=1}^{\tilde{p}}\frac{\partial F}{\partial I_{j}}\underset{=X_{\boldsymbol{\theta}}(I_{j})}{\underbrace{\left(\sum_{\ell=1}^{p}\chi_{\ell}\frac{\partial I_{j}}{\partial\theta_{\ell}}\right)}}\nonumber\\
    &=\sum_{j=1}^{\tilde{p}}\frac{\partial F}{\partial I_{j}}\underset{=0}{\underbrace{X_{\boldsymbol{\theta}}(I_{j})}}=0,
\end{align}
which is the desired result.
\end{proof}

\end{document}
